\documentclass[10pt]{article}

\usepackage[a4paper,margin=0.72in]{geometry}
\usepackage{amsmath}
\usepackage[T1]{fontenc}

\setlength{\parindent}{0pt}
\setlength{\parskip}{0.42em}

\begin{document}

\begin{center}
{\Large Kerr Linear Inversion in \texttt{bayRing}}
\end{center}

\textbf{Kerr-only model.}
For a Kerr ringdown without tail terms or quadratic amplitudes, and with the
non-amplitude Kerr quantities fixed, the complex strain is linear in the QNM
amplitudes:
\begin{equation}
  h(t_j) = h_{\rm fixed}(t_j)
  + \sum_{k=1}^{N} C_k\,q_k(t_j),
  \qquad
  C_k = A_k e^{i\phi_k}.
\end{equation}
Here \(q_k(t)\) is the Kerr basis waveform for mode \(k=(\ell,m,n)\), including
the fixed Kerr frequency and damping time. The term \(h_{\rm fixed}\) contains
any constant offset or already-fixed mode amplitudes. In the simplest
\texttt{QNM-modes = 220} example, the only unknown is
\(C_{220}=A_{220}e^{i\phi_{220}}\).

Because the model is linear in the complex amplitudes \(C_k\), we do not need to
do stochastic inference or nonlinear minimization in this case. Once the
non-amplitude quantities are fixed, the amplitudes follow from one deterministic
linear-algebra solve.

\textbf{Real weighted system.}
The code solves for real and imaginary amplitude parts,
\begin{equation}
  x =
  \left(\operatorname{Re}C_1,\operatorname{Im}C_1,\ldots,
        \operatorname{Re}C_N,\operatorname{Im}C_N\right)^T .
\end{equation}
After subtracting \(h_{\rm fixed}\), the real and imaginary data components are
stacked and divided by their corresponding errors. The same weighting is applied
to the basis columns, giving a real design matrix \(B\) and data vector \(d\):
\begin{equation}
  d \simeq Bx .
\end{equation}
For each Kerr mode, \(B\) has two columns: the weighted waveform generated by
\(C_k=1\) and by \(C_k=i\), where \(i\) is the complex unit, \(i^2=-1\).

\textbf{Closed-form solve.}
The weighted least-squares normal equations are
\begin{equation}
  F x = r,
  \qquad
  F = B^T B,
  \qquad
  r = B^T d .
\end{equation}
The implementation symmetrizes \(F\), diagonalizes it, and floors small
eigenvalues using the configured tolerance \(\epsilon\):
\begin{equation}
  F = V\Lambda V^T,
  \qquad
  \hat{x}
  =
  V\,\operatorname{diag}\!\left[
      \frac{1}{\max(\lambda_i,\epsilon)}
    \right]V^T r .
\end{equation}
Each recovered complex amplitude
\(\hat{C}_k=\hat{x}_{2k-1}+i\hat{x}_{2k}\) is then written back in the
\texttt{bayRing} parameterization,
\begin{equation}
  \ln A_k = \log |\hat{C}_k|,
  \qquad
  \phi_k = \arg(\hat{C}_k) \bmod 2\pi .
\end{equation}
The corresponding one-sigma error estimate is obtained from the same weighted
Fisher matrix. In Cartesian amplitude coordinates the covariance is
\begin{equation}
  \Sigma_x =
  V\,\operatorname{diag}\!\left[
      \frac{1}{\max(\lambda_i,\epsilon)}
    \right]V^T .
\end{equation}
The code then propagates \(\Sigma_x\) to the \((\ln A_k,\phi_k)\)
parameterization using the local Jacobian of the polar transformation.

\textbf{Configuration.}
Use \texttt{method = Linear-inversion}. For the Kerr-only setup discussed here,
leave each \(\ln A\), \(\phi\) amplitude pair free, and fix or omit all
nonlinear parameters. The eigenvalue floor is set by the configured linear
inversion tolerance. A minimal example is
\begin{center}
\small
\texttt{config\_files/config\_SXS\_0305\_Kerr\_220\_}\\
\texttt{linear\_inversion.ini}
\end{center}

\end{document}
